\section*{Chapitre \ref{ch:b1:acetals-protection} --- Les composés carbonylés~: additions nucléophiles, acétals et protection}

\begin{solution}{exo:b1:acetals-protection:1}
\ce{CH3CH(OH)OCH3}~: \omterm{def:b1:acetals-protection:hemiacetal}{hémiacétal}. \ce{CH3CH(OCH3)2}~: \omterm{def:b1:acetals-protection:hemiacetal}{acétal}.
\ce{CH3CH(OH)2}~: hydrate. \ce{CH3OCH3}~: éther. \ce{(CH3)2C(OCH2CH3)2}~:
\omterm{def:b1:acetals-protection:hemiacetal}{acétal} (d'une cétone).
\end{solution}

\begin{solution}{exo:b1:acetals-protection:2}
\ce{CH3CH2CHO + 2CH3OH <=> CH3CH2CH(OCH3)2 + H2O}~; la propanone donne
avec l'éthane-1,2-diol le 2,2-diméthyl-1,3-dioxolane et de l'eau,
\ce{(CH3)2CO + HOCH2CH2OH <=> C5H10O2 + H2O}.
\end{solution}

\begin{solution}{exo:b1:acetals-protection:3}
L'hydroxyde de sodium, le bromure de méthylmagnésium et le borohydrure
de sodium le laissent inchangé~; l'acide sulfurique dilué en solution
aqueuse l'hydrolyse.
\end{solution}

\begin{solution}{exo:b1:acetals-protection:4}
Le OH du carbone 4 s'additionne sur le carbone 1 de l'aldéhyde~: un
cycle à cinq atomes, quatre carbones et un oxygène, avec un OH sur le
carbone voisin de l'oxygène du cycle.
\end{solution}

\begin{solution}{exo:b1:acetals-protection:5}
Protonation de C=O~; addition du méthanol sur le carbone~; perte de
\ce{H+}~: \omterm{def:b1:acetals-protection:hemiacetal}{hémiacétal}~; protonation de son OH~; perte d'eau conduisant à
\ce{CH3CH=O+CH3}~; addition d'une seconde molécule de méthanol~; perte
de \ce{H+}~: \ce{CH3CH(OCH3)2}. Les deux protons captés sont restitués~:
l'acide est un \omterm{def:b1:elementary-steps:catalyst}{catalyseur}.
\end{solution}

\begin{solution}{exo:b1:acetals-protection:6}
Protonation d'un groupe OCH3~; perte de méthanol conduisant à
\ce{(CH3)2C=O+CH3}~; addition de l'eau~;
perte de \ce{H+}~: \omterm{def:b1:acetals-protection:hemiacetal}{hémiacétal}~; protonation, perte de méthanol, perte de
\ce{H+}~: propanone. L'excès d'eau maintient $Q$ inférieur à $K$ pour
l'hydrolyse, qui est totale.
\end{solution}

\begin{solution}{exo:b1:acetals-protection:7}
Une molécule d'eau par \omterm{def:b1:acetals-protection:hemiacetal}{acétal}~: $0.250 \times 18.0 = \qty{4.5}{g}$, soit
environ \qty{4.5}{mL}. Au bout d'une heure, $3.2/18.0 = \qty{0.178}{mol}$~:
taux de conversion de \qty{71}{\percent}.
\end{solution}

\begin{solution}{exo:b1:acetals-protection:8}
Avec le diol, une molécule en remplace deux~: la réaction ne diminue pas
le nombre de molécules (un aldéhyde et un diol donnent un \omterm{def:b1:acetals-protection:hemiacetal}{acétal} et une
molécule d'eau), et le second OH est maintenu près du carbone qui
réagit, ce qui ferme facilement le cycle. Ces deux effets rendent
l'\omterm{def:b1:acetals-protection:hemiacetal}{acétal} cyclique plus favorable.
\end{solution}

\begin{solution}{exo:b1:acetals-protection:9}
L'acide détruirait l'\omterm{def:b1:organomagnesium:organometallic}{organomagnésien} (transfert de proton)~; il faut
aussi éliminer les traces d'eau et de diol.
\end{solution}

\begin{solution}{exo:b1:acetals-protection:10}
1. Protéger la cétone de la 4-chlorobutan-2-one sous forme d'\omterm{def:b1:acetals-protection:hemiacetal}{acétal}
cyclique (éthane-1,2-diol, \omterm{def:b1:elementary-steps:catalyst}{catalyseur} acide, Dean--Stark). 2. Magnésium
dans l'éther anhydre~: l'\omterm{def:b1:organomagnesium:organometallic}{organomagnésien} du chlorure protégé, qui n'est
plus détruit par son propre carbonyle. 3. Ajouter la propanone~;
hydrolyse par une solution aqueuse de chlorure d'ammonium~: l'alcool
tertiaire. 4. Solution aqueuse acide diluée~: \omterm{def:b1:acetals-protection:protecting-group}{déprotection} en
5-hydroxy-5-méthylhexan-2-one.
\end{solution}

\begin{solution}{exo:b1:acetals-protection:11}
$Q = [\text{acétal}][\ce{H2O}]/([\text{aldéhyde}][\ce{CH3OH}]^2)$~: un
large excès de méthanol, élevé au carré au dénominateur, maintient
$Q < K$ jusqu'à ce qu'il reste peu d'aldéhyde. L'élimination de l'eau
agit sur le numérateur~; employer le méthanol comme solvant est plus
simple mais exige un produit facile à séparer.
\end{solution}

\begin{solution}{exo:b1:acetals-protection:12}
Seul le OH de l'\omterm{def:b1:acetals-protection:hemiacetal}{hémiacétal} peut partir sous forme d'eau en donnant un
cation stabilisé par l'oxygène voisin du cycle (\ce{C=O+})~; les autres
groupes OH devraient partir de carbones ordinaires, en donnant des
cations non stabilisés. Le méthanol s'additionne sur le cation stabilisé~:
un \omterm{def:b1:acetals-protection:hemiacetal}{acétal} méthylique.
\end{solution}

\begin{solution}{pb:b1:acetals-protection:1}
\textbf{1.} \ce{BrMgCH2CH2CH2CHO}.
\textbf{2.} Son carbone de type \omterm{def:b1:electronic-effects:intermediates}{carbanion} attaquerait l'aldéhyde d'une
autre molécule (ou le sien)~: \ce{RMgBr} s'additionne sur
$\mathrm{R'CHO}$, en donnant un \omterm{def:b1:alcohol-activation:alkoxide}{alcoolate}.
\textbf{3.} L'aldéhyde, contre l'\omterm{def:b1:organomagnesium:organometallic}{organomagnésien}.
\textbf{4.} Les \omterm{def:b1:acetals-protection:hemiacetal}{acétals} sont stables vis-à-vis des \omterm{def:b1:organomagnesium:organometallic}{organomagnésiens} et
des bases, et sont éliminés par une solution aqueuse acide.
\textbf{5.} Il se forme plus complètement (un diol pour deux molécules
d'alcool) et s'hydrolyse facilement.
\textbf{6.} $15.09/150.9 = \qty{0.1000}{mol}$.
\textbf{7.} $1.20 \times 0.1000 \times 62.0 = \qty{7.44}{g}$.
\textbf{8.} \ce{BrCH2CH2CH2CHO + HOCH2CH2OH <=> C6H11BrO2 + H2O}
(2-(3-bromopropyl)-1,3-dioxolane).
\textbf{9.} Il active le carbonyle et permet la perte d'eau lors de la
seconde étape~; il est régénéré.
\textbf{10.} Le toluène et l'eau distillent ensemble, se condensent et
tombent dans le collecteur~; l'eau coule au fond et y reste, le toluène
déborde et retourne dans le ballon.
\textbf{11.} $Q$ contient $[\ce{H2O}]$ à son numérateur~; éliminer l'eau
maintient $Q < K$, et la réaction progresse.
\textbf{12.} L'eau est plus dense et non \omterm{def:b1:intermolecular-forces-solvents:miscibility}{miscible} au toluène~: elle
s'accumule au fond~; seule la \omterm{def:b1:extent-q-and-k:system}{phase} supérieure atteint le bras de retour.
\textbf{13.} Quand l'eau cesse de s'accumuler, au volume attendu.
\textbf{14.} $0.1000 \times 194.9 = \qty{19.5}{g}$.
\textbf{15.} \ce{C6H11BrO2 + Mg -> C6H11BrMgO2} (éther anhydre).
\textbf{16.} Le réactif s'additionne sur \ce{(CH3)2CO}~: l'\omterm{def:b1:alcohol-activation:alkoxide}{alcoolate}
\ce{(CH3)2C(O^-)CH2CH2CH2-} sur la chaîne protégée, avec \ce{MgBr+}.
\textbf{17.} L'acide retirerait trop tôt le \omterm{def:b1:acetals-protection:protecting-group}{groupe protecteur} et
pourrait déshydrater l'alcool tertiaire.
\textbf{18.} $0.0700 \times 174.0 = \qty{12.2}{g}$.
\textbf{19.} L'\omterm{def:b1:acetals-protection:hemiacetal}{acétal} + \ce{H2O} $\to$ l'aldéhyde + éthane-1,2-diol
(\omterm{def:b1:elementary-steps:catalyst}{catalyseur} acide).
\textbf{20.} Sa \omterm{def:b1:alcohol-activation:dehydration}{déshydratation} exige un \omterm{def:b1:predominant-reaction:strong-acid}{acide fort} et un chauffage~; une
solution aqueuse faiblement acide et froide hydrolyse l'\omterm{def:b1:acetals-protection:hemiacetal}{acétal} bien plus
vite.
\textbf{21.} Le OH du carbone 5 s'additionne sur le carbone 1 de
l'aldéhyde~: un cycle à six atomes (cinq carbones et un oxygène) portant
deux groupes méthyle sur le carbone voisin de l'oxygène du cycle et un
OH sur l'autre carbone voisin de cet oxygène.
\textbf{22.} $0.0700 \times 0.90 = \qty{0.0630}{mol}$, soit \qty{8.19}{g}.
\textbf{23.} $0.0630/0.1000 = \qty{63}{\percent}$ (\omterm{def:b1:acetals-protection:protecting-group}{protection} supposée
totale).
\textbf{24.} $0.1000 \times 18.0 = \qty{1.80}{g}$.
\textbf{25.} $1.80/0.997 =$ \textbf{\qty{1.8}{mL} d'eau} dans le collecteur.
\end{solution}
